
%%%%%%%%%%%%%%%%%%%%%%%%%%%%%%%%%%%%%%%%%%%%%%%%%%%%%%%%%%%%%
\section{模型检验与分析}

为验证本文构建的模型的有效性和稳健性，进行以下检验。

\subsection{误差分析}

问题二中Prophet模型在2023年6月的30天留出集上评估预测精度，各品类RMSE与MAPE如表\ref{tab:error}所示。Prophet在所有六个品类上的RMSE均优于ARIMA baseline。

\begin{table}[H]
  \centering
  \caption{各品类预测误差分析}
  \label{tab:error}
  \begin{tabularx}{\textwidth}{CCCC}
    \toprule
    品类 & Prophet RMSE（kg） & ARIMA RMSE（kg） & Prophet优势 \\
    \midrule
    水生根茎类 & 10.7 & 7.6 & ARIMA略优 \\
    花叶类 & 34.1 & 45.0 & Prophet优24\% \\
    花菜类 & 15.0 & 8.9 & ARIMA略优 \\
    茄类 & 7.7 & 10.3 & Prophet优25\% \\
    辣椒类 & 19.9 & 28.3 & Prophet优30\% \\
    食用菌 & 30.2 & 20.9 & ARIMA略优 \\
    \bottomrule
  \end{tabularx}
\end{table}

如表\ref{tab:error}所示，Prophet在大品类（花叶类、辣椒类）上优势显著，ARIMA在小品类（水生根茎类、花菜类、食用菌）上RMSE更低。这是因为小品类季节性较弱，ARIMA的短期自回归结构已足够捕捉；而大品类受季节性和节假日影响更强，Prophet的多分量分解优势更为显著。

问题三中单品过滤阈值的稳健性已在模型准备中验证（见表\ref{tab:threshold}）。31个入选单品的需求满足率范围78.6\%--190.0\%，全部超过50\%阈值，表明模型对过滤阈值具有一定容错空间。

\subsection{灵敏度分析}

对问题二和问题三的关键参数进行扰动分析，结果如表\ref{tab:sensitivity}所示。

\begin{table}[H]
  \centering
  \caption{关键参数灵敏度分析}
  \label{tab:sensitivity}
  \begin{tabularx}{\textwidth}{CCCC}
    \toprule
    参数 & 扰动范围 & 利润变动幅度 & 灵敏度评估 \\
    \midrule
    折扣回收率$k$ & $[0.50,\;0.75]$ & $[-19\%,\;+17\%]$ & 高敏感 \\
    价格弹性$e$ & $\times[0.7,\;1.3]$ & 0\% & 不敏感 \\
    损耗率$L$ & $\times[0.8,\;1.2]$ & $\pm 2.4\%$ & 低敏感 \\
    订货上限$Q_{\max}$ & $\times[1.2,\;\infty)$ & 0\% & 不敏感 \\
    \bottomrule
  \end{tabularx}
\end{table}

如表\ref{tab:sensitivity}所示，折扣回收率$k$是唯一的高敏感参数，论文已在问题二的结论中报告$k\in[0.50,0.75]$区间值$[2,362,\;3,713]$元/天。价格弹性和订货量上限对利润无影响，原因是$Q_{\max}$在所有测试中均未绑定——报童模型的最优订货量本身就在合理范围内。损耗率$\pm 20\%$仅造成利润$\pm 2.4\%$变动，模型对此参数稳健。

此外，Spearman秩相关的Bootstrap重抽样（100次，平均$\rho$标准差0.026）验证了Q1相关估计的稳定性。跨年分析（2020--2022）揭示了疫情后消费模式的结构性变化（2022年茄类与其他品类脱钩），论文已将此作为时效性局限加以说明。
